\chapter{REVIEW OF LITERATURE}
\label{ch:literature}

\section{Introduction}
\label{sec:lit_intro}

This chapter provides a critical synthesis of the literature relevant
to the study.  It is structured around three converging streams of
scholarship: (i) the \emph{theoretical} foundations from
communications, media-management and cultural-industries literature
which inform the conceptual framework outlined in
Chapter~\ref{ch:introduction}; (ii) the \emph{thematic} body of work
on the streaming turn, OTT platforms, audience behaviour and the
Indian media-and-entertainment industry; and (iii) the \emph{contextual}
literature on Tamil cinema specifically and on the COVID-19 disruption
of cultural industries.  The review identifies converging arguments,
contested findings and -- most importantly for this study -- gaps that
the present thesis addresses.  More than three hundred sources were
read in the preparation of this chapter; one hundred and fifty-plus
are cited in the present synthesis and in the chapters that follow.

The chapter is organised as follows.  Section~\ref{sec:theory}
examines four theoretical traditions that are jointly operationalised
in the conceptual framework.  Section~\ref{sec:streaming_turn} surveys
the global streaming-studies literature.  Section~\ref{sec:india_me}
traces the pre-pandemic trajectory of the Indian
media-and-entertainment sector.  Section~\ref{sec:ott_india} reviews
work on Indian and South-Asian OTT.  Section~\ref{sec:audience}
synthesises cinema-audience studies.
Section~\ref{sec:tamilcinema} covers the dedicated Tamil-cinema
literature.  Section~\ref{sec:covid_cultural} surveys
COVID-19-and-cultural-industries scholarship.
Section~\ref{sec:windowing} addresses windowing theory and
industrial-reconfiguration scholarship.  Section~\ref{sec:starsystem}
examines the star-versus-content literature.
Section~\ref{sec:gap} consolidates the identified gaps.
Section~\ref{sec:lit_summary} summarises.

\section{Theoretical Foundations}
\label{sec:theory}

\subsection{Diffusion of Innovations Theory}
\label{sec:doi}

Rogers's \citeyearpar{rogers2003} \emph{Diffusion of Innovations}
remains the dominant theory for explaining the rate and pattern of
adoption of a new technology, product or practice in a social system.
The five attributes of an innovation -- relative advantage,
compatibility, complexity, trialability and observability -- and the
S-shaped adoption curve formed by innovators, early adopters, early
majority, late majority and laggards have been adapted to the study
of new media adoption \citep{lin2003,vishwanath2005,straub2009}.

In the context of streaming, several scholars have employed
diffusion frameworks.  \citet{cha2013} examined the adoption of OTT
in the United States and found relative advantage and compatibility
to be the strongest predictors.  \citet{shin2018} extended the
framework to South-Korean OTT, integrating perceived enjoyment as a
context-specific construct.  \citet{shukla2022} applied the model to
Indian OTT during COVID-19 and found a sharp compression of the
diffusion timeline -- what would normally take five to seven years
of adoption being achieved in eighteen months.  \citet{jaiswal2022}
identified ``forced trialability'' under lockdown as an explanatory
variable absent from Rogers's original schema.

For our purposes, three implications of Diffusion theory matter.
First, the adoption rate observed in our viewer survey
(Chapter~\ref{ch:viewership}) needs to be benchmarked against the
S-curve to determine whether the post-pandemic equilibrium is one of
saturation, of plateauing, or of continuing growth.  Second, the
moderating role of demographics -- age, education, income, urban /
rural residence -- is a standard component of diffusion analysis and
maps directly onto our research questions.  Third, the
``compatibility'' attribute has a particular bite in a linguistically
distinctive market: the lack of \emph{Tamil-language} originals on
several global platforms in 2018--2019 was a compatibility frictionit
that subsequent localisation has reduced \citep{punathambekar2013,
mehta2020streaming}.

\subsection{Uses and Gratifications Theory}
\label{sec:ug}

The Uses and Gratifications (U\&G) tradition reverses the deterministic
``what does media do to people'' framing of early effects research and
asks instead ``what do people do with media'' \citep{katz1973,
blumler1979}.  Audiences are conceptualised as goal-directed actors
selecting media to satisfy specific cognitive, affective, integrative
and tension-release needs.

The framework has been extended to the internet and streaming era by
\citet{ruggiero2000}, \citet{papacharissi2000}, \citet{sundar2013} and
others.  In the OTT-specific literature, four broad gratification
clusters have been recurrently identified:

\begin{enumerate}[itemsep=2pt]
  \item \textbf{Content-seeking} -- desire for original, niche or
        previously unavailable content \citep{sundet2021,jenner2018,
        lotz2017}.
  \item \textbf{Convenience} -- portability, time-shifting,
        device-flexibility \citep{lotz2017,tryon2013}.
  \item \textbf{Social} -- co-viewing, conversation around content,
        social-media participation \citep{pittman2016,lim2020}.
  \item \textbf{Hedonic} -- relaxation, mood-management, escape
        \citep{tefertiller2017,reer2017}.
\end{enumerate}

Indian-context studies have largely confirmed these clusters
\citep{singh2021,gupta2022uses,parmar2021} while adding a culturally
specific gratification: language- and culture-affinity, particularly
salient for regional-language viewers \citep{punathambekar2013,
nakassis2017}.  Our Chapter~\ref{ch:viewership} analysis tests the
relative weight of these gratifications in the post-pandemic Tamil
viewer profile.

\subsection{Disruptive Innovation Theory}
\label{sec:disruption}

Christensen's \citeyearpar{christensen1997} theory of disruptive
innovation -- subsequently elaborated by Christensen, Raynor \&
McDonald \citep{christensen2015} -- argues that incumbent firms are
typically vulnerable to entrants who initially serve under-served or
non-consumer market segments with simpler, cheaper offerings, and who
then move up-market.  Theorists have debated whether streaming is
``truly'' disruptive of theatrical exhibition or merely
\emph{adjacent} \citep{wayne2018,fagerjord2018}.  \citet{lobato2019}
argues that the dichotomy is misleading: streaming and theatre are
better understood as occupying overlapping but differentiated
windows in a single value chain.

For the present study, the key analytical task is to determine
empirically whether, in Tamil cinema specifically, the post-pandemic
relationship between OTT and theatre is one of substitution
(disruption) or complementarity (co-existence).
Chapter~\ref{ch:viewership} provides direct evidence on the
question.

\subsection{Cultural Industries Paradigm}
\label{sec:cultural_industries}

The cultural-industries tradition, traceable to \citet{adorno1944},
\citet{garnham1990} and developed in the contemporary form by
\citet{hesmondhalgh2019}, treats cultural production as a hybrid
economic-aesthetic activity governed by a distinctive
``risk-management'' logic.  Five characteristics are typically
emphasised: high fixed cost / low marginal cost economics; high
demand uncertainty (``nobody knows''); reliance on stars and
franchises as risk-reduction devices; a complex division of creative
and reproductive labour; and the political-economic significance of
cultural ``texts'' beyond their immediate market value
\citep{caves2000,deuze2007,deyandolf2014}.

\citet{caldwell2008} extends the paradigm to ``production-cultures''
research, focusing on how industry workers themselves theorise and
narrate change.  \citet{curtin2014} introduces the concept of
``media capital'' -- spatially clustered, scale-economy-bearing,
trajectory-dependent agglomerations.  Both inform our reading of
Chennai as a media-capital site and our specialist-interview design
in Chapters~\ref{ch:methodology} and~\ref{ch:industry_reconfig}.

\subsection{Convergence and Platform Studies}
\label{sec:convergence}

\citet{jenkins2006} defines media convergence as the flow of content
across platforms, the cooperation among multiple media industries
and the migratory behaviour of audiences.  Subsequent
platform-studies scholarship \citep{plantin2018,gillespie2018,
nieborg2018} has shifted attention to the infrastructural, economic
and discursive role of platforms as ``programmable'' intermediaries
\citep{vandijck2018}.  \citet{lobato2019} -- in his foundational
\emph{Netflix Nations} -- offers the most comprehensive treatment of
streaming as a transnational, jurisdictionally fragmented and
linguistically negotiated phenomenon.  These platform-studies
insights inform our analysis of platform preference (Chapter~5,
Figure~5.8) and platform-driven content commissioning
(Chapter~6).

\section{The Global Streaming Turn}
\label{sec:streaming_turn}

The 2010s witnessed what \citet{lotz2017,lotz2022} terms the ``second
phase'' of internet-distributed television: the maturation of Netflix
from a US DVD-rental incumbent into the dominant global
direct-to-consumer streaming platform; the entry of Amazon Prime
Video, Disney+, HBO Max, Apple TV+ and a host of regional services;
and the transformation of the production economics of long-form
audiovisual content.  Several recurring findings in this literature
are pertinent.

\textit{First}, the global streaming wars have been characterised by
sustained capital intensity and content-spend escalation
\citep{lobato2019,fitzgerald2019,jenner2018}, with Netflix alone
spending over USD 17 billion on content in 2022
\citep{economist2022netflix}.

\textit{Second}, despite global ambition, regional and local content
has emerged as a critical battleground.  \citet{lotz2022}, drawing on
work by \citet{straubhaar2007cultural} on cultural proximity, argues
that local-language originals retain a structural advantage even in
markets penetrated by global platforms.

\textit{Third}, the streaming wars have produced a re-windowing
debate.  \citet{cunningham2019} review the historical ``windows''
system -- theatrical, home video, premium TV, free TV -- and argue
that the post-2018 window restructuring is the most significant
re-engineering of cinema distribution since the introduction of
home video in 1976.

\textit{Fourth}, scholars have begun to write the global streaming
history with explicit attention to non-Western formations
\citep{punathambekar2013,jin2019,kanozia2021}.

For our specific purpose, the global streaming literature provides
two analytical anchors: the historical baseline of windowing
behaviour, and the international comparators against which Tamil
specifics can be evaluated.  Several of these comparators reappear
in Chapter~7's discussion section.

\section{The Indian Media-and-Entertainment Sector: Pre-Pandemic
Trajectory}
\label{sec:india_me}

The annual FICCI--EY ``M\&E Outlook'' reports \citep{ficciey2018,
ficciey2019,ficciey2020} provide the most comprehensive longitudinal
data on the Indian sector.  Key trajectory features prior to 2020:

\begin{itemize}[itemsep=2pt]
  \item Television was the largest sub-sector by revenue (\textsterling 78,7
        thousand crore in 2019), with subscription accounting for
        over half of revenue \citep{ficciey2020}.
  \item Film industry revenue stood at \textsterling 19,1 thousand crore
        in 2019, with theatrical accounting for 60--62\% of the
        total \citep{ficciey2020}.
  \item Digital media (which includes OTT) reached \textsterling 22,1
        thousand crore in 2019, growing at a CAGR of 26\% over
        2016--2019 \citep{ficciey2020}.
  \item Print and radio were in stagnation or decline.
\end{itemize}

\citet{kohlikhandekar2013}'s \emph{The Indian Media Business} -- now
in its fourth and fifth editions -- remains the canonical
single-volume historical and structural reference for the Indian
media business.  Pre-2020 OTT specifically is treated by
\citet{punathambekar2013}, \citet{mehta2020streaming},
\citet{kanozia2021} and the trade reports of
\citet{kpmg2019,bcg2018india}.

A consistent theme is the rising salience of regional-language
content within both broadcast television and OTT, with Tamil and
Telugu identified as among the fastest-growing regional sub-segments
\citep{ficciey2020,bcg2018india,statista2019}.  This is the
trajectory which the COVID-19 pandemic interrupted and accelerated
in non-uniform ways.

\section{OTT in India and South Asia}
\label{sec:ott_india}

The Indian OTT market grew from approximately \textsterling 4,400 crore in
2018 to \textsterling 7,800 crore in 2020 and is projected to exceed
\textsterling 21,000 crore by 2026 \citep{ficciey2024,kpmg2023,statista2024}.
The Indian streaming literature can be organised into four sub-clusters:

\subsection{Industry-structure studies}

\citet{kanozia2021}, \citet{singh2022ottstructure} and
\citet{parmar2021} map the entry, consolidation and competitive
dynamics of Indian OTT platforms.  \citet{economist2022jiocinema}
and \citet{economictimes2024jiomerger} document the role of Reliance
Jio in commoditising mobile data, in launching JioCinema in 2022 and
in driving the 2024 Disney--Reliance merger.

\subsection{Content and programming studies}

\citet{mehta2019} and \citet{punathambekar2018} examine the
emergence of an Indian original-content aesthetic that departs from
both Bollywood norms and global Netflix formulas.
\citet{nayak2021} examines the regulatory contestation around the
Information Technology (Intermediary Guidelines and Digital Media
Ethics Code) Rules, 2021.

\subsection{Audience studies}

\citet{singh2021,sharma2020otta,gupta2022uses,parmar2021,
chakraborty2022} are among the dozens of audience studies of Indian
OTT, typically using survey methods and either U\&G or technology
acceptance frameworks.  Most are limited by national or pan-Indian
sampling that obscures regional-language specificities.

\subsection{Regional-language streaming studies}

A small but growing literature addresses regional-language streaming
specifically.  \citet{nakassis2020} examines Tamil cinema's flirtation
with web series; \citet{velayutham2019} discusses the early Tamil
streaming services; \citet{aha2022launchreport} (a trade-association
white paper) describes the launch of Aha Tamil and its first two
years.  Notably, no peer-reviewed monograph has yet treated Tamil OTT
as a coherent subject.  The present study aims to fill this gap.

\section{Cinema Audience Studies and Viewership Research}
\label{sec:audience}

Cinema audience studies have a long methodological history, ranging
from the Payne Fund Studies of the 1930s \citep{jowett1996} through
post-war effects research, the British cultural-studies turn
\citep{stacey1994,kuhn2002} and the digital-era audience studies of
\citet{barker2008}, \citet{tryon2013} and \citet{ross2008}.  In the
Indian context, classic field studies include
\citet{dickey1993,prasad1998,vasudevan2010,mazumdar2007}, and
contemporary work by \citet{ganti2012,nakassis2017,gopalan2002,
mukherjee2020}.

For audience-survey methodology, the standard references are
\citet{babbie2020}, \citet{groves2009} and \citet{tourangeau2000} on
survey-research design, and \citet{taylorbogdan2015}, \citet{
strauss1998} and \citet{saldana2021} on qualitative data analysis.
Methodological combinations of survey, interview and trade-data
analysis appear in \citet{redvall2013}, \citet{noh2020} and
\citet{singh2022ottstructure}.

\citet{ormax2022,ormax2023,ormax2024} -- though industry rather than
academic -- provide the most current systematic data on Indian
cinema-going behaviour with breakdowns by language, age and city
tier.  These reports are extensively used in
Chapter~\ref{ch:industry_overview} as a sanity-check on our primary
viewer survey.

\section{Studies on the Tamil Film Industry}
\label{sec:tamilcinema}

The dedicated literature on the Tamil film industry can be grouped
chronologically and thematically:

\subsection{Historical scholarship}

\citet{baskaran1996,baskaran2009} provide the canonical historical
treatments, tracing Tamil cinema from the silent era through the
talkies of the 1930s, the post-war Dravidianist phase, the
M.\,G.\,Ramachandran star-political compact of the 1960s--1970s, the
Sivaji-MGR-Rajini lineage and the contemporary multistar economy.
\citet{pandian1992} provides a political-cultural reading of the MGR
phenomenon.  \citet{hardgrave1973} is the early classic on the
politics-cinema continuum.

\subsection{Linguistic-cultural studies}

\citet{ramaswamy1997} examines Tamil-language ideology and its
intersections with cinema.  \citet{nakassis2010,nakassis2017}
provides ethnographic accounts of Tamil cinema's distinctive
star-system, signage practice and audience-address.

\subsection{Industrial-economic studies}

\citet{velayutham2008} edited an early multi-author volume on Tamil
cinema as a cultural industry; \citet{krishnan2009} analyses the
political economy of distribution; \citet{ravindran2018} examines
the late-2010s budget escalation; \citet{kohlikhandekar2013}
provides industry-wide data within an all-India frame.

\subsection{Genre, gender and representation studies}

\citet{rajan2002,velayutham2008,nakassis2017,prasad2014,
sundar2017tamil} examine Tamil cinema's genre repertoires, gender
representation and aesthetic conventions.  These works do not
directly address the streaming era, but they establish the
cultural-textual baseline against which our genre-shift findings in
Chapter~\ref{ch:viewership} are interpreted.

\subsection{The streaming-era gap}

There is, to date, no peer-reviewed monograph that systematically
documents the post-2020 transformation of Tamil cinema, although
trade journalism has produced rich descriptive material
\citep{filmcompanion2021directott,filmcompanion2022mid,
behindwoods2023,galatta2023}.  Several master's dissertations and
non-refereed conference papers have begun to address the topic
\citep{krishnan2022mphil,raja2023conference,sundaresan2023case},
but a peer-reviewed empirical study of the kind attempted in this
thesis appears not to exist.

\section{COVID-19 and the Cultural Industries}
\label{sec:covid_cultural}

A substantial international literature on the COVID-19 disruption of
cultural industries emerged in 2020--2023 \citep{flew2021,florida2020,
comunian2021,banks2020,khlystova2022}.  Key findings include:

\begin{itemize}[itemsep=2pt]
  \item Cultural workers were among the most acutely affected
        labour categories, with disproportionate impact on
        freelance, project-based and live-performance workers
        \citep{comunian2021,banks2020}.
  \item The shock accelerated digitalisation in production,
        distribution and consumption \citep{flew2021,khlystova2022}.
  \item Recovery has been highly uneven, with some sub-sectors
        (e.g., gaming, streaming) benefiting while others
        (e.g., live-performance, cinema-exhibition) lagging
        \citep{ficciey2023,unesco2022culture}.
  \item National policy responses have varied substantially in
        scale, design and effectiveness
        \citep{flew2021,unesco2021culture}.
\end{itemize}

For India, \citet{ficciey2021,kpmg2021,kpmg2022,bcg2022} document
sectoral impact and recovery.  Cinema-specific COVID literature
includes \citet{vasileva2021}, \citet{nichols2021},
\citet{tryon2021covid} and the special issue of \emph{Studies in
Documentary Film} on COVID-era documentary
\citep{aufderheide2021special}.  No published peer-reviewed study
to date isolates the Tamil-language film industry's COVID experience
in detail.

\section{Industrial Reconfiguration and Windowing Theory}
\label{sec:windowing}

The economics of release windows -- theatrical, home video, pay-TV,
free-TV -- has been examined by \citet{vogel2020}, \citet{epstein2010},
\citet{waterman2005} and \citet{dewenter2017}.  \citet{cunningham2019}
review the implications of streaming for the windows system.
\citet{vogel2020} models the trade-off between revenue maximisation
and consumer-welfare in window-length decisions.  COVID-era
re-windowing has been studied by \citet{schauerte2021},
\citet{waterman2022pandemic} and \citet{garcia2023streaming}.  The
Tamil-specific specialist data presented in
Chapter~\ref{ch:industry_reconfig} adds to this empirical base.

Industrial-reconfiguration scholarship more broadly draws on
organisational-economics traditions \citep{teece2018,helfat2018},
co-evolutionary theory \citep{murmann2003} and dynamic-capabilities
arguments \citep{teece2018,helfat2018}.  These frameworks are useful
in interpreting the speed and pattern of Tamil-industry adaptation,
which we discuss in Section~\ref{sec:windowing_findings} of
Chapter~\ref{ch:industry_reconfig}.

\section{The Star System and the Star-vs-Content Debate}
\label{sec:starsystem}

The star-system literature is large and conceptually rich.  Classic
references include \citet{dyer1979,dyer1986}, \citet{morin1960},
\citet{deCordova1990} and \citet{mcdonald2013star}.  Indian
star-studies include \citet{prasad2014},
\citet{majumdar2009starsindia}, \citet{rajadhyaksha2009},
\citet{nakassis2017} and \citet{punathambekar2013}.

A particularly contested question post-2020 has been whether
streaming has \emph{weakened} the star system in favour of
content-driven decision-making, or whether it has \emph{relocated}
star power into a new, multi-platform celebrity economy
\citep{cunningham2019,sundet2021,lotz2022}.  Empirical evidence on
the weakening hypothesis appears in
\citet{schauerte2021,garcia2023streaming}, while
\citet{lotz2022,nakassis2017} caution against premature obituary of
star-driven cultural industries.  In Tamil cinema specifically the
question has explosive relevance because of the unusually intense
star-fan-club apparatus described in
Section~\ref{sec:tamilcharacteristics} and tested empirically in
Chapter~\ref{ch:viewership}.

\section{The Indian Streaming Wars: A Timeline, 2018--2025}
\label{sec:streaming_wars}

To anchor the contextual literature in a coherent temporal narrative,
this section reconstructs a brief but substantive timeline of the
Indian streaming wars over the eight years 2018--2025, drawing on
\citet{ficciey2018,ficciey2019,ficciey2020,ficciey2021,ficciey2022,
ficciey2023,ficciey2024}, \citet{kpmg2019,kpmg2021,kpmg2022,kpmg2023},
\citet{bcg2018india,bcg2022}, \citet{statista2019,statista2023,
statista2024}, \citet{economist2022jiocinema,economictimes2024jiomerger,
economist2022netflix} and a curated archive of trade-press articles.

\textit{2018 -- the early plateau}.  Hotstar (then majority-owned by
Star India / Twenty-First Century Fox) was the largest Indian SVOD by
subscribers, supported by IPL cricket rights.  Netflix India had
fewer than a million paying subscribers; Amazon Prime Video had
launched in 2016 and was investing aggressively in Indian originals.
Sun NXT (2017 launch) was the largest Tamil-first SVOD with a
catalogue rather than originals strategy.  Voot (2016 launch by
Viacom18) and ZEE5 (2018 launch) were nascent.  Subscription revenue
across the Indian SVOD market in 2018 totalled an estimated INR
2{,}100 crore \citep{ficciey2019}.

\textit{2019 -- IPL inflection and regional kick-off}.  The Hotstar
Premium subscription gained scale on the back of cricket; Disney
acquired Twenty-First Century Fox globally, bringing Hotstar into
the Disney portfolio.  The first Indian Netflix originals to attract
critical attention -- \emph{Sacred Games} (2018), \emph{Delhi Crime}
(2019) -- demonstrated the format proposition.  Aha was announced
late in the year, with launch planned for early 2020 in Telugu and
later in Tamil.  Tamil-language originals were minimal.

\textit{2020 -- the pandemic surge}.  COVID-19 produced the structural
shock that defines this thesis.  Direct-to-OTT releases of completed
films -- \emph{Soorarai Pottru} (Suriya), \emph{Penguin}, \emph{Vakeel
Saab} -- created a precedent for Tamil and Telugu cinema.
Subscriptions accelerated sharply.  Aha launched in Tamil in late
2020 / early 2021 \citep{aha2022launchreport}.  Disney+ launched as a
brand in India, integrated with Hotstar.

\textit{2021 -- platform proliferation and acquisition-price spike}.
Acquisition prices for Tamil OTT rights briefly more than doubled.
Trade-press attention to direct-to-OTT economics intensified
\citep{filmcompanion2021directott,variety2021directott}.  The
2021 IT Rules (Information Technology (Intermediary Guidelines and
Digital Media Ethics Code) Rules, 2021) introduced new content-classification
and grievance-redressal obligations; \citet{nayak2021} provides a
detailed legal analysis.  By end-2021 the Indian SVOD market had
crossed 95 million paying subscribers \citep{kpmg2022,ficciey2022}.

\textit{2022 -- consolidation and content-spend re-discipline}.
Several global platforms reviewed their India content-spend in line
with global cost-discipline imperatives \citep{economist2022netflix}.
Hotstar lost the IPL streaming rights to Reliance / Viacom18's
JioCinema in mid-2022, an inflection point for the Indian streaming
market \citep{economist2022jiocinema}.  Acquisition-price multiples
began correcting.  Theatrical recovery for Tamil cinema reached a
post-pandemic high (Figure~\ref{fig:bo_trend}).

\textit{2023 -- the regional-original turn}.  Tamil-language originals
on Aha Tamil, Sun NXT, Disney+ Hotstar, Amazon Prime Video, ZEE5 and
JioCinema collectively crossed 100 titles for the first time.
Acquisition-price corrections continued.  The trade press began to
discuss the ``mid-budget renaissance''
\citep{filmcompanion2022mid,behindwoods2023,galatta2023}.

\textit{2024 -- the Disney--Reliance merger}.  Walt Disney Company
and Reliance Industries announced the merger of their Indian media
operations -- Star India / Disney+ Hotstar with Viacom18 / JioCinema
-- to create the dominant Indian streaming entity by paying
subscribers \citep{economictimes2024jiomerger,thehindu2024merger}.
The merger has consequences for content commissioning, advertiser
pricing and competition policy that are still unfolding at the time
of writing.

\textit{2025 -- the Tamil-language equilibrium}.  The Tamil-language
SVOD market has, by mid-2025, settled into an oligopoly of Amazon
Prime Video, Sun NXT, JioHotstar (the merged entity), Aha Tamil and
Netflix, with YouTube and the AVOD tier of JioHotstar continuing to
provide a substantial free-streaming long-tail.  Subscriber
penetration in Tamil-Nadu households has reached approximately 60\%
(Figure~\ref{fig:ott_penetration}).  Acquisition prices have
stabilised above the 2018 baseline at the star tier and below the
2018 baseline at the small / indie tier (Figure~\ref{fig:ott_price}),
indicating a market that is structurally larger but more disciplined
than at the 2020--21 peak.

This timeline is the immediate macro-economic context against which
the present thesis's specialist-narrated reconfiguration must be
read.  No single platform dominates the Tamil market in the way
Netflix dominates several Western markets, and the regional-language
SVOD market is more fragmented than the Hindi-language market.

\section{Identified Research Gaps}
\label{sec:gap}

Synthesising the foregoing, we identify seven gaps which the present
thesis addresses:

\begin{description}[style=nextline,leftmargin=2.4em]
\item[Gap 1.] No peer-reviewed empirical study has yet documented
       Tamil-cinema viewership transformation post-COVID with a
       primary, demographically stratified, sub-state representative
       sample.  The 400-respondent viewer survey of
       Chapter~\ref{ch:viewership} addresses this gap.
\item[Gap 2.] No study has yet integrated demand-side viewership
       evidence with supply-side specialist evidence for Tamil
       cinema in a single mixed-methods design.  The combined
       analysis of Chapters~\ref{ch:viewership}--\ref{ch:industry_reconfig}
       addresses this gap.
\item[Gap 3.] Existing diffusion-of-innovation studies of Indian
       OTT do not explicitly isolate the Tamil-Nadu /
       Tamil-language sub-population.  The age-, location- and
       income-stratified analyses of Chapter~\ref{ch:viewership}
       provide this missing isolation.
\item[Gap 4.] The post-2020 evolution of the theatrical-window
       structure in Tamil cinema has been treated only descriptively
       in trade journalism; no academic work has yet quantified the
       distribution across budget tiers using specialist data.
       Chapter~\ref{ch:industry_reconfig} addresses this gap.
\item[Gap 5.] No published quantitative work has documented the
       OTT-acquisition-price index trajectory across film tiers
       specific to Tamil cinema.  Section~\ref{sec:price_index} of
       Chapter~\ref{ch:industry_reconfig} addresses this gap.
\item[Gap 6.] The star-versus-content debate has not been
       empirically tested for Tamil cinema with a probability-style
       survey.  Section~\ref{sec:star_content_results} of
       Chapter~\ref{ch:viewership} provides direct evidence.
\item[Gap 7.] No integrated framework has been proposed to explain
       the post-pandemic Tamil-cinema equilibrium.  The
       Hybrid-Equilibrium framework of Chapter~\ref{ch:conclusion}
       offers one.
\end{description}

\section{Chapter Summary}
\label{sec:lit_summary}

This chapter has reviewed and synthesised the literature relevant
to the study, working through four theoretical traditions
(Diffusion of Innovations, Uses \& Gratifications, Disruptive
Innovation and Cultural Industries), nine thematic clusters (global
streaming, Indian M\&E, Indian OTT, audience studies, Tamil-cinema
historical, Tamil-cinema industrial, COVID-cultural,
windowing-theory, star-system) and the dedicated Tamil-cinema
literature.  Seven research gaps were identified and aligned with
specific chapters of the thesis.  The next chapter sets out the
research methodology by which these gaps are addressed empirically.
